\documentclass[11pt]{article}

% arXiv standard packages
\usepackage[utf8]{inputenc}
\usepackage[T1]{fontenc}
\usepackage{hyperref}
\usepackage{url}
\usepackage{booktabs}
\usepackage{amsfonts}
\usepackage{amsmath}
\usepackage{nicefrac}
\usepackage{microtype}
\usepackage{fancyhdr}
\usepackage{graphicx}
\usepackage{titlesec}
\usepackage{geometry}
\usepackage{enumitem}
\usepackage{listings}
\usepackage{xcolor}
\usepackage{natbib}
\usepackage{setspace}
\usepackage{CJKutf8}

% Page geometry
\geometry{a4paper, margin=1in}
\onehalfspacing

% Hyperlink setup
\hypersetup{
    colorlinks=true,
    linkcolor=blue,
    filecolor=magenta,
    urlcolor=cyan,
    citecolor=blue,
}

% Code listing style
\definecolor{codegreen}{rgb}{0,0.6,0}
\definecolor{codegray}{rgb}{0.5,0.5,0.5}
\definecolor{codepurple}{rgb}{0.58,0,0.82}
\definecolor{backcolour}{rgb}{0.95,0.95,0.92}

\lstdefinestyle{mystyle}{
    backgroundcolor=\color{backcolour},
    commentstyle=\color{codegreen},
    keywordstyle=\color{magenta},
    numberstyle=\tiny\color{codegray},
    stringstyle=\color{codepurple},
    basicstyle=\ttfamily\footnotesize,
    breakatwhitespace=false,
    breaklines=true,
    captionpos=b,
    keepspaces=true,
    numbers=left,
    numbersep=5pt,
    showspaces=false,
    showstringspaces=false,
    showtabs=false,
    tabsize=2
}
\lstset{style=mystyle}

% Title formatting
\titleformat{\section}{\large\bfseries}{\thesection}{1em}{}
\titleformat{\subsection}{\normalsize\bfseries}{\thesubsection}{1em}{}

% Header
\pagestyle{fancy}
\fancyhf{}
\fancyhead[L]{\small\textit{Preprint --- not peer reviewed}}
\fancyhead[R]{\small\thepage}
\renewcommand{\headrulewidth}{0.4pt}

\begin{document}

\begin{CJK}{UTF8}{gbsn}

\title{\textbf{CittaVerse Narrative Scorer v0.5:\\Six-Dimension Assessment for Chinese\\Autobiographical Memory Quality}}

\author{
    \textbf{CittaVerse Team} \\
    \begin{CJK}{UTF8}{gbsn}一念万相科技 (CittaVerse)\end{CJK} \\
    Hangzhou, China \\
    \texttt{cittaverse@gmail.com}
}

\date{March 2026}

\maketitle

\end{CJK}

\begin{abstract}
\noindent Reminiscence therapy (RT) has demonstrated moderate effects on cognitive function and psychological well-being in older adults with mild cognitive impairment (MCI). However, scalable delivery and objective quality assessment remain significant challenges. Current practice relies on manual narrative coding, which is labor-intensive, subjective, and impractical for large-scale deployment. We present the CittaVerse Narrative Scorer v0.5, an automated assessment tool that evaluates Chinese autobiographical memories across six dimensions: event richness, temporal coherence, causal coherence, emotional depth, identity integration, and information density distribution. Our neuro-symbolic approach combines rule-based feature extraction with research-derived scoring heuristics, achieving full automation without requiring large language model APIs. The scorer is deployed in an ongoing randomized controlled pilot study (N=50, 2-week intervention) and provides immediate, interpretable feedback to both participants and clinicians. This technical report describes the methodology, algorithm design, and evaluation framework. The source code is openly available under the MIT license to support reproducibility and community validation. We position this work as a foundational contribution to AI-enhanced digital therapeutics for cognitive health in aging populations.

\vspace{0.5em}
\noindent\textbf{Keywords}: reminiscence therapy, narrative assessment, autobiographical memory, Chinese NLP, digital therapeutics, mild cognitive impairment, older adults, neuro-symbolic AI
\end{abstract}

%% ============================================================
%% SECTION 1: INTRODUCTION
%% ============================================================

\section{Introduction}

\subsection{Background and Motivation}

The global population is aging rapidly. By 2050, the number of people aged 60 years and older is projected to reach 2.1 billion, with China alone accounting for over 400 million \citep{un2024}. Alongside demographic shifts comes a growing prevalence of cognitive impairment. Mild cognitive impairment (MCI) affects 10--20\% of adults over 65 and represents a critical window for intervention before progression to dementia \citep{petersen2016}.

Reminiscence therapy (RT)---the structured recall and discussion of past experiences---has emerged as a promising non-pharmacological intervention for older adults with MCI. Meta-analyses demonstrate moderate effects on cognitive function (Hedges' $g = 0.35$--$0.45$), depression ($g = 0.40$), and quality of life \citep{pu2025, ni2026}. Digital RT platforms extend these benefits through scalable, home-based delivery, reducing barriers of geography, mobility, and therapist availability \citep{wang2026}.

However, a critical gap persists: \textbf{how do we objectively assess the quality of narratives produced during RT?} Current practice relies on manual coding using established protocols such as the Autobiographical Interview (AI) \citep{levine2002}, which distinguishes internal (episodic) from external (semantic) details. While psychometrically validated, manual coding requires trained raters, achieves only moderate inter-rater reliability ($\kappa = 0.6$--$0.7$), and cannot provide real-time feedback \citep{spencer2020}. This bottleneck limits both clinical scalability and research throughput.

\subsection{Problem Statement}

The challenge of automated narrative assessment in Chinese older adults with MCI involves three intersecting constraints:

\begin{enumerate}[label=(\arabic*)]
    \item \textbf{Linguistic specificity}: Chinese narrative structure differs from English in temporal marking, causal connectives, and self-reference patterns \citep{chen2021}. Tools developed for Western languages cannot be directly transferred.
    \item \textbf{Clinical validity}: Scoring dimensions must align with established cognitive constructs (e.g., episodic specificity, coherence) while remaining sensitive to MCI-related changes \citep{murphy2023}.
    \item \textbf{Computational efficiency}: Deployment in resource-constrained settings (community centers, home use) requires lightweight algorithms that do not depend on cloud-based large language models (LLMs).
\end{enumerate}

\subsection{Our Contribution}

We present the CittaVerse Narrative Scorer v0.5, addressing these challenges through:

\begin{enumerate}[label=(\arabic*)]
    \item \textbf{Six-dimension framework}: Extending the classic internal/external dichotomy \citep{levine2002} with temporal coherence, causal coherence, emotional depth, identity integration, and information density distribution---dimensions identified through systematic review of 2024--2026 evidence on narrative and cognition in aging \citep{shankar2025, yang2026, kensinger2026}.
    \item \textbf{Chinese-language optimization}: Vocabulary lists for temporal markers, causal connectives, self-references, and emotion words curated from Chinese psycholinguistic resources and validated against native speaker intuition.
    \item \textbf{Neuro-symbolic design}: Rule-based feature extraction (symbolic) combined with research-derived scoring heuristics (neural-inspired), achieving interpretability without sacrificing automation.
    \item \textbf{Open-source implementation}: Complete source code, test suite, and example narratives released under MIT license to support community validation and extension.
\end{enumerate}

\subsection{Paper Organization}

Section~\ref{sec:related} reviews related work in narrative assessment, LLM-based evaluation, and technology acceptance for older adults. Section~\ref{sec:methodology} describes the six-dimension framework and theoretical foundations. Section~\ref{sec:implementation} details the algorithm design and implementation. Section~\ref{sec:validation} outlines the validation strategy and pilot evaluation framework. Section~\ref{sec:ethics} discusses ethical considerations, limitations, and future directions. Section~\ref{sec:conclusion} concludes.

%% ============================================================
%% SECTION 2: RELATED WORK
%% ============================================================

\section{Related Work}\label{sec:related}

\subsection{Autobiographical Memory and Reminiscence Therapy}

The Autobiographical Interview (AI) \citep{levine2002} remains the gold standard for assessing episodic specificity in narrative recall. The AI distinguishes:
\begin{itemize}
    \item \textbf{Internal details}: Episodic elements tied to a specific time and place
    \item \textbf{External details}: Semantic information, repetitions, and off-topic content
\end{itemize}

Research consistently shows that older adults with MCI produce fewer internal details and more external details compared to healthy controls \citep{irish2024, barnhofer2025}. However, the AI requires 20--30 minutes of rater training and 10--15 minutes per narrative for coding---impractical for scalable deployment.

Recent meta-analyses confirm RT efficacy but highlight assessment heterogeneity as a limitation. \citet{pu2025} reviewed 23 digital RT studies and found wide variation in outcome measures (cognitive tests, mood scales, qualitative interviews), complicating cross-study comparison. \citet{ni2026} emphasized the need for standardized, objective narrative quality metrics.

\subsection{Automated Narrative Assessment}

Natural language processing (NLP) offers promising alternatives to manual coding. Early work used surface features (word count, sentence length, type-token ratio) to discriminate MCI from healthy controls with 70--80\% accuracy \citep{fraser2024}. More recent approaches leverage:

\begin{itemize}
    \item \textbf{Latent Semantic Analysis (LSA)}: Captures semantic coherence through word co-occurrence matrices \citep{thomas2025}
    \item \textbf{Syntactic complexity metrics}: Clause density, embedding depth, dependency length \citep{lu2025}
    \item \textbf{Topic modeling}: LDA-based extraction of thematic structure \citep{zhang2025}
\end{itemize}

However, these methods were developed primarily for English and focus on dementia detection rather than therapeutic progress monitoring.

\subsection{LLM-Based Evaluation}

The emergence of large language models (LLMs) has transformed automated text evaluation. ``LLM-as-a-Judge'' approaches achieve human-level agreement on tasks ranging from essay scoring \citep{zheng2025} to dialogue quality \citep{liu2025}. In healthcare, LLMs have been applied to:
\begin{itemize}
    \item Mental health symptom extraction from clinical notes \citep{chen2025}
    \item Suicide risk assessment from social media posts \citep{wang2025}
    \item Therapy session quality evaluation \citep{goldberg2025}
\end{itemize}

A 2026 RCT by \citet{limbic2026} demonstrated that LLM-personalized interventions significantly increased engagement and improved outcomes in anxiety/depression treatment. However, LLM-based approaches raise concerns about:
\begin{itemize}
    \item \textbf{Bias}: Training data underrepresents older adults and non-Western languages \citep{bender2021}
    \item \textbf{Opacity}: Black-box scoring limits clinical interpretability
    \item \textbf{Cost}: API-based deployment incurs ongoing expenses
\end{itemize}

\subsection{Technology Acceptance for Older Adults}

Successful deployment of digital therapeutics requires attention to technology acceptance. The Senior Technology Acceptance Model (STAM) \citep{mitzner2024} identifies key constructs: technology anxiety, privacy concerns, perceived usefulness, and facilitating conditions.

Extended TAM studies (2025--2026) in Chinese older adults highlight additional factors:
\begin{itemize}
    \item \textbf{Emotional safety}: Comfort with AI discussing personal memories \citep{zhang2025}
    \item \textbf{Perceived personalization}: Belief that the system adapts to individual needs \citep{liu2025b}
    \item \textbf{Social influence}: Family/caregiver endorsement \citep{chen2026}
\end{itemize}

\subsection{Neuro-Symbolic AI for Healthcare}

Neuro-symbolic approaches combine neural networks' pattern recognition with symbolic reasoning's interpretability \citep{garcez2025}. In healthcare applications, this paradigm has been applied to clinical decision support with explainable rules, medical question answering with knowledge graph grounding, and drug interaction checking with logic constraints. Our work extends this paradigm to narrative assessment, using rule-based feature extraction (symbolic) with research-derived scoring heuristics (neural-inspired).

\subsection{Gap Analysis}

Table~\ref{tab:gap} summarizes the landscape. The CittaVerse Narrative Scorer addresses the intersection of Chinese language, automated scoring, clinical interpretability, and open-source availability---a combination not previously available.

\begin{table}[h]
\centering
\caption{Landscape of narrative assessment tools}
\label{tab:gap}
\begin{tabular}{lcccc}
\toprule
\textbf{Tool} & \textbf{Language} & \textbf{Automation} & \textbf{Dimensions} & \textbf{Open Source} \\
\midrule
Autobiographical Interview & English & Manual & 2 & No \\
LSA-based methods & English & Automatic & 1 & Some \\
LLM-as-a-Judge & Multi & Automatic & Variable & No \\
\textbf{CittaVerse Scorer (ours)} & \textbf{Chinese} & \textbf{Automatic} & \textbf{6} & \textbf{Yes} \\
\bottomrule
\end{tabular}
\end{table}

%% ============================================================
%% SECTION 3: METHODOLOGY
%% ============================================================

\section{Methodology}\label{sec:methodology}

\subsection{Six-Dimension Framework}

Our framework extends the classic internal/external dichotomy \citep{levine2002} with five additional dimensions identified through systematic review:

\begin{table}[h]
\centering
\caption{Six-Dimension Assessment Framework}
\label{tab:dimensions}
\begin{tabular}{lll}
\toprule
\textbf{Dimension} & \textbf{Construct} & \textbf{Rationale} \\
\midrule
Event Richness & Episodic specificity & Core AI internal detail count \\
Temporal Coherence & Temporal grounding & Time markers anchor events \\
Causal Coherence & Narrative causality & Causal connectives indicate meaning-making \\
Emotional Depth & Affective engagement & Emotion words correlate with consolidation \\
Identity Integration & Self-referential processing & Self-pronouns indicate autobiographical reasoning \\
Information Density & Central vs.\ peripheral & Emotional arousal enhances central details \\
\bottomrule
\end{tabular}
\end{table}

\subsection{Theoretical Foundations}

\subsubsection{Event Richness and Episodic Specificity}

The Autobiographical Interview defines internal details as ``episodic elements tied to a specific time and place'' \citep{levine2002}. Research shows that older adults with MCI produce significantly fewer internal details ($d = 0.6$--$0.8$) compared to healthy controls \citep{irish2024, barnhofer2025}. Our event extraction algorithm identifies:
\begin{itemize}
    \item \textbf{Action events}: Verbs indicating specific actions
    \item \textbf{Perceptual events}: Sensory details
    \item \textbf{Emotional events}: Affective experiences
\end{itemize}

\subsubsection{Temporal and Causal Coherence}

Temporal coherence reflects the ability to situate events within a coherent timeline. Chinese narrative structure uses explicit temporal markers more frequently than English \citep{chen2021, li2025}. Causal coherence indicates meaning-making---the process of connecting events into a coherent life story \citep{habermas2020, mcadams2024}. Both dimensions are impaired in MCI and dementia \citep{murphy2023, devito2025}.

\subsubsection{Emotional Depth and Identity Integration}

Emotional depth captures the affective richness of narratives. Research shows that emotional arousal enhances memory consolidation through amygdala-hippocampal interactions \citep{kensinger2026, reed2024}. Identity integration reflects self-referential processing---the extent to which narratives engage with the narrator's sense of self \citep{conway2023, adler2025}. Autobiographical reasoning (connecting past events to current identity) is a key mechanism of RT efficacy \citep{westerhof2024}.

\subsubsection{Information Density Distribution}

Recent work by \citet{kensinger2026} and \citet{picard2025} demonstrates that emotional arousal enhances central information encoding at the expense of peripheral details. This ``tunnel memory'' effect is mediated by amygdala filtering. Our scorer computes the central/peripheral ratio, with an optimal target of 60/40 based on empirical evidence.

\subsection{Scoring Algorithm}

For each dimension, we compute a density-based score normalized to 0--100:
\begin{equation}
\text{score}_i = \min\left(\frac{\text{raw\_count}_i}{\text{text\_length}} \times \text{norm\_factor}_i \times 100,\; 100\right)
\end{equation}

Normalization factors are derived from pilot data (N=50 Chinese older adults):

\begin{table}[h]
\centering
\caption{Normalization Factors by Dimension}
\label{tab:normalization}
\begin{tabular}{llc}
\toprule
\textbf{Dimension} & \textbf{Factor} & \textbf{Weight ($w_i$)} \\
\midrule
Event Richness & 400 & 0.15 \\
Temporal Coherence & 200 & 0.15 \\
Causal Coherence & 150 & 0.15 \\
Emotional Depth & 100 & 0.20 \\
Identity Integration & 50 & 0.15 \\
Information Density & 100 & 0.20 \\
\bottomrule
\end{tabular}
\end{table}

The composite score is a \textbf{weighted} mean:
\begin{equation}
\text{Composite} = \sum_{i=1}^{6} w_i \times \text{score}_i
\end{equation}

Emotional Depth and Information Density receive higher weights ($w = 0.20$) based on their stronger association with therapeutic outcomes in reminiscence therapy \citep{westerhof2024, kensinger2026}. Letter grades are assigned: S (90--100), A (80--89), B (70--79), C (60--69), D (50--59), F ($<$50).

%% ============================================================
%% SECTION 4: IMPLEMENTATION
%% ============================================================

\section{Implementation}\label{sec:implementation}

\subsection{Architecture Overview}

The scorer is implemented in Python 3.8+ with zero external dependencies (standard library only). This design choice enables:
\begin{itemize}
    \item \textbf{Offline deployment}: No internet connection required
    \item \textbf{Low resource usage}: $<$50MB RAM, suitable for community center computers
    \item \textbf{Easy maintenance}: No dependency conflicts or version drift
\end{itemize}

\begin{lstlisting}[language=bash, caption={Narrative Scorer Repository Structure}]
narrative-scorer/
  README.md              # Usage documentation
  LICENSE                # MIT License
  requirements.txt       # Empty (no external deps)
  src/
    scorer.py            # Core scoring logic (14.5KB)
  examples/
    sample_input.txt     # Example Chinese narrative
    sample_output.json   # Example scoring output
  tests/
    test_scorer.py       # 11 unit tests
\end{lstlisting}

\subsection{Feature Extraction}

\subsubsection{Vocabulary Lists}

We curate four vocabulary lists from Chinese psycholinguistic resources:

\begin{CJK}{UTF8}{gbsn}
\textbf{Temporal Markers (24 items)}: 然后，接着，随后，之后，后来，之前，当时，现在，过去，曾经，已经，正在，将要，首先，其次，最后，同时，于是，便，才，就，再，又，还

\textbf{Causal Connectives (13 items)}: 因为，所以，因此，因而，于是，从而，致使，导致，由于，为此，既然，可见，总之

\textbf{Self-References (6 items)}: 我，我的，我自己，咱，咱们，俺

\textbf{Emotion Words (32 items)}: 开心，快乐，高兴，难过，悲伤，痛苦，愤怒，生气，害怕，恐惧，惊讶，惊奇，平静，安宁，焦虑，紧张，兴奋，激动，失望，沮丧，满足，幸福，孤独，寂寞，温暖，感动，自豪，骄傲，后悔，遗憾，感激，感恩
\end{CJK}

Total vocabulary: 75 items. Future versions will expand to 200+ based on corpus analysis.

\subsubsection{Event Extraction}

Events are extracted using regex patterns for Chinese verb structures:
\begin{itemize}
    \item \textbf{Action verbs}: V + \begin{CJK}{UTF8}{gbsn}了\end{CJK} / V + \begin{CJK}{UTF8}{gbsn}过\end{CJK} / V + \begin{CJK}{UTF8}{gbsn}着\end{CJK}
    \item \textbf{Perceptual verbs}: \begin{CJK}{UTF8}{gbsn}看到，听到，闻到，感到，觉得\end{CJK}
    \item \textbf{Motion verbs}: \begin{CJK}{UTF8}{gbsn}去，来，到，走，跑\end{CJK}
\end{itemize}

Event boundaries are detected using temporal markers and sentence boundaries.

\subsection{Performance Characteristics}

\begin{table}[h]
\centering
\caption{Performance Metrics}
\label{tab:performance}
\begin{tabular}{ll}
\toprule
\textbf{Metric} & \textbf{Value} \\
\midrule
Latency (100 chars) & $\sim$5ms \\
Latency (1000 chars) & $\sim$15ms \\
Throughput & $\sim$60 narratives/second \\
Memory usage & $<$50MB \\
Test coverage & 11 unit tests, 100\% pass \\
\bottomrule
\end{tabular}
\end{table}

\subsection{Output Format}

JSON output includes six dimension scores (0--100), composite score with letter grade (S/A/B/C/D/F), natural language feedback in Chinese, and raw feature counts for debugging and research transparency.

%% ============================================================
%% SECTION 5: VALIDATION STRATEGY
%% ============================================================

\section{Validation Strategy}\label{sec:validation}

\subsection{Pilot RCT Design}

The scorer is deployed in an ongoing randomized controlled pilot study:

\begin{table}[h]
\centering
\caption{Pilot RCT Parameters}
\label{tab:rct}
\begin{tabular}{ll}
\toprule
\textbf{Parameter} & \textbf{Value} \\
\midrule
Sample size & N=50 (25 intervention, 25 control) \\
Duration & 2 weeks \\
Intervention & Daily RT sessions with AI guidance \\
Control & Waitlist \\
Primary outcome & MoCA change \\
Secondary outcomes & GDS (depression), narrative quality scores \\
Assessment points & Baseline, Week 1, Week 2, Follow-up (4 weeks) \\
\bottomrule
\end{tabular}
\end{table}

\subsection{Validation Metrics}

We will assess:
\begin{enumerate}[label=(\arabic*)]
    \item \textbf{Convergent validity}: Correlation with manual AI coding ($r > 0.7$ expected)
    \item \textbf{Discriminant validity}: MCI vs.\ healthy control separation ($d > 0.5$ expected)
    \item \textbf{Test-retest reliability}: ICC $> 0.8$ over 1-week interval
    \item \textbf{Sensitivity to change}: Pre-post intervention effect size ($d > 0.3$ expected)
    \item \textbf{Clinical utility}: Therapist satisfaction survey (Likert 1--5, target $>4.0$)
\end{enumerate}

\subsection{Power Analysis}

For the pilot RCT (N=50), we have 80\% power to detect $d = 0.8$ (large effect) and 60\% power to detect $d = 0.5$ (medium effect). A full-scale RCT (N=200, planned Q3 2026) will have 80\% power to detect $d = 0.4$ (medium-small effect).

%% ============================================================
%% SECTION 6: ETHICAL CONSIDERATIONS
%% ============================================================

\section{Ethical Considerations}\label{sec:ethics}

\subsection{Privacy and Data Security}

Narrative data contains sensitive personal information. We implement local storage (no cloud upload by default), optional encryption for data at rest, explicit consent for research use, and right to deletion at any time.

\subsection{Algorithmic Fairness}

Our vocabulary lists are curated from standard Mandarin resources. Limitations include limited dialect coverage (currently Mandarin only), education bias, and cultural specificity based on Han Chinese norms. Future work will expand to dialect support and cross-cultural validation.

\subsection{Clinical Boundaries}

The scorer is a \textbf{research and wellness tool}, not a diagnostic device. We explicitly do not claim to diagnose MCI or dementia, do not replace clinical assessment, and recommend professional consultation for concerning results.

\subsection{Emotional Safety}

Recalling past memories can evoke strong emotions. We provide pre-session warnings about potential emotional responses, option to skip difficult topics, crisis resource links (suicide hotline, mental health services), and therapist oversight in clinical deployments.

%% ============================================================
%% SECTION 7: LIMITATIONS AND FUTURE WORK
%% ============================================================

\section{Limitations and Future Work}

\subsection{Current Limitations}

\begin{enumerate}[label=(\arabic*)]
    \item Rule-based event extraction may miss implicit events or metaphorical language
    \item Limited vocabulary (75 items); expansion needed for broader coverage
    \item Mandarin-only; no dialect support yet
    \item No LLM enhancement; pure rule-based approach may miss nuanced patterns
    \item Pilot validation only; large-scale RCT needed for definitive evidence
\end{enumerate}

\subsection{Future Directions}

\begin{enumerate}[label=(\arabic*)]
    \item \textbf{LLM-enhanced scoring}: Hybrid approach using LLMs for event extraction while retaining interpretable dimension scores
    \item \textbf{Vocabulary expansion}: 200+ markers driven by corpus analysis
    \item \textbf{Dialect support}: Cantonese, Hokkien, Shanghainese
    \item \textbf{Multimodal integration}: Photo/audio cues for richer prompting
    \item \textbf{Longitudinal tracking}: Narrative trajectories over months/years
    \item \textbf{Cross-cultural adaptation}: English, Spanish, Japanese versions
    \item \textbf{ASR integration}: Speech-to-text pipeline for end-to-end scoring
\end{enumerate}

%% ============================================================
%% SECTION 8: CONCLUSION
%% ============================================================

\section{Conclusion}\label{sec:conclusion}

We present the CittaVerse Narrative Scorer v0.5, the first automated assessment tool for Chinese autobiographical memory quality. Our six-dimension framework extends classic narrative assessment with temporal coherence, causal coherence, emotional depth, identity integration, and information density distribution. The neuro-symbolic design achieves full automation without LLM API dependencies, enabling offline deployment in resource-constrained settings.

The scorer is deployed in an ongoing pilot RCT (N=50) and will undergo comprehensive validation including convergent validity with manual coding, discriminant validity between MCI and healthy controls, and sensitivity to therapeutic change. We release the source code under MIT license to support community validation and extension.

As global demographics shift toward older populations, scalable, evidence-based interventions for cognitive health become increasingly critical. Automated narrative assessment represents one piece of this puzzle---enabling objective, real-time feedback that supports both therapeutic progress and research advancement.

%% ============================================================
%% ACKNOWLEDGMENTS
%% ============================================================

\begin{CJK}{UTF8}{gbsn}
\section*{Acknowledgments}
We thank the CittaVerse team for their contributions to this work. This research is supported by 一念万相科技 (CittaVerse).
\end{CJK}

%% ============================================================
%% REFERENCES
%% ============================================================

\bibliographystyle{plainnat}
\bibliography{references}

\end{document}
